\section{Titans and Praetorians}

These powerful engines of destruction cannot follow the same rules as infantry or tanks. The specific rules governing them are described below.

\subsection{Purchasing Titans and Praetorians}

Most Titans can be equipped with a variety of weapons suited to their mission. The available chassis, weapon descriptions, and army points costs are listed in the relevant Codexes.

The procedure for equipping a Titan is always the same. First purchase the chassis, then purchase its weapons where required. Every non-optional weapon mount must be equipped with a weapon.

A weapon that may be mounted anywhere can be purchased for either an arm mount or a carapace mount.

The same assault weapon cannot be selected twice for a single Titan. Assault weapons are marked with ``(Assault)'' after their name.

\subsection{Movement}

Titans and Praetorians are cumbersome machines with limited manoeuvrability. They move in straight lines and may turn by no more than a total of 90\textdegree{} during their movement. This turning allowance may be divided between several separate turns.

They may turn up to 90\textdegree{} during the Movement Phase and may also turn up to 90\textdegree{} when making an assault Consolidation move, if necessary.

Titans and Praetorians may move in reverse like any other unit. Each centimetre moved in reverse counts as 2~cm of movement.

\subsection{Overwatch Fire}

The weapons of Titans and Praetorians may perform Overwatch Fire independently of one another. A Titan or Praetorian may therefore perform Overwatch Fire with one of its weapons and subsequently activate its remaining weapons during the Combat Phase.

When this happens, reveal its First Fire order. It cannot perform any further Overwatch Fire during the current turn.

When performing Overwatch Fire against a detachment, you must declare every weapon that will fire at it before rolling any dice.

\subsection{Damage Resolution}

Damage affects Titans and Praetorians differently from other units. When a Titan or Praetorian is hit, follow this procedure:

\begin{enumerate}
    \item \textbf{Saving Throw:} Make saving throws normally.

    \item \textbf{Hit Location:} Hits may strike different locations, such as the hull, head, or a weapon. Consult the Titan's hit-location chart, using either the front or rear chart depending on whether the attack came from within its front or rear arc. Roll 1D6 to determine which location was hit.

    \item \textbf{Consult the Damage Table:} For each failed saving throw, the Titan or Praetorian loses one Wound and suffers damage to the location that was hit.

    Damage is progressive. Each unsaved hit inflicts the first damage result that has not yet been suffered by that location. For example, if a Titan has already suffered the first damage result to one of its legs, the next hit to that leg inflicts the second damage result, and so forth.
\end{enumerate}

For weapons, there are as many separate hit locations as there are weapon entries on the Titan's profile. The attacker chooses which weapon is hit, even if that weapon has already been destroyed by previous damage.

\subsection{Damage Effects}

Some damage results have additional effects. For example:

\begin{description}
    \item[Weapon Damaged:] The weapon cannot be used until it is repaired.

    \item[Weapon Destroyed:] The weapon is permanently destroyed and can no longer be used.

    \item[Immobilised:] The Titan becomes immobilised. See Immobilised Base.
\end{description}

\subsection{Assault}

Assaults involving Titans and Praetorians are resolved in the same manner as assaults involving other bases.

If a Titan or Praetorian loses a duel, it suffers one hit with no saving throw permitted. The hit is then resolved using the same procedure as a shooting attack.

\subsection{Destruction}

A Titan is destroyed if it is reduced to zero Wounds. Depending on the damage it has suffered, its destruction may cause additional effects:

\begin{description}
    \item[Reactor Explosion:] The Titan or Praetorian is destroyed. Every base within 3D6~cm suffers a hit on a roll of 2+ with AP~$-1$, treated as a shooting attack originating from the Titan or Praetorian.

    Bases inside the exploding Titan or Praetorian cannot make an Emergency Exit as a result of this damage.

    The explosion also affects garrisoned bases if at least half of their building is covered by the explosion.

    \item[Titan Fall:] Roll a scatter die to determine the direction in which the Titan falls. To determine which bases are hit, place a 12~cm template in contact with the Titan's base in the direction of the fall.

    If a ``Hit'' is rolled, the Titan collapses upon itself with no additional effect.

    Bases beneath the template suffer a hit on a roll of 4+ with AP~$-3$. Buildings suffer one automatic hit with the \textit{Damages Buildings (AP $-4$/3)} ability.
\end{description}

\subsection{Repairs}

If a damage result is marked ``Damaged'', it may be repaired during the End-of-Turn Effects step of the End Phase.

Only the most recent damage result suffered by a location can be repaired. Any previous damage results affecting that location are permanent.

Roll 1D6. On a result of 4+, the damage is repaired and its additional effects are cancelled.

Only the effect of the damage is repaired. For the purpose of determining future damage results, the location is still considered to have suffered that damage.
